% Generated draft: requires mathematical and source review.
\documentclass[11pt]{article}
\usepackage[margin=25mm]{geometry}
\usepackage{fontspec}
\setmainfont{Latin Modern Roman}
\usepackage{amsmath,amssymb,mathtools,unicode-math}
\setmathfont{Latin Modern Math}
\usepackage{enumitem,needspace,xurl,hyperref}
\setlength{\parindent}{0pt}
\setlength{\parskip}{5pt}
\setlength{\emergencystretch}{3em}
\newcommand{\N}{\mathbb{N}}
\newcommand{\Z}{\mathbb{Z}}
\newcommand{\Q}{\mathbb{Q}}
\newcommand{\R}{\mathbb{R}}
\newcommand{\C}{\mathbb{C}}
\newcommand{\F}{\mathbb{F}}
\newcommand{\A}{\mathbb{A}}
\newcommand{\PP}{\mathbb{P}}
\newcommand{\E}{\mathbb{E}}
\newcommand{\eps}{\varepsilon}
\newcommand{\dd}{\,\mathrm{d}}
\begin{document}
\section*{UnsolvedMath ID: 30000668}
\textbf{Source problem code:} catalogue entry
\subsection{Definitions and mathematical statement}
Fix the ground field $k$ intended by the source. Let $G$ be a connected reductive algebraic $k$-group and let $\mathfrak g=\operatorname{Lie}(G)$. Use conjugation on $G$ and the adjoint action on $\mathfrak g$. A dominant rational map $f:G\dashrightarrow\mathfrak g$ is $G$-equivariant if
\[
f(hgh^{-1})=\operatorname{Ad}(h)f(g)
\]
as an identity of rational maps. Since $\dim G=\dim\mathfrak g$, such a dominant map is generically finite; its degree is the function-field degree
\[
\deg(f)=[k(G):f^*k(\mathfrak g)].
\]
When this class of maps is nonempty, the source defines
\[
\operatorname{Cay}(G)=\min_f\deg(f).
\]
The questions are:
\begin{enumerate}
\item Does there exist a connected reductive algebraic group in the intended field setting with $\operatorname{Cay}(G)>2$?
\item Does there exist such a group that is simple?
\end{enumerate}
Here the first scope follows the class for which the source defines Cayley degree; the subsequent phrase ``an algebraic group'' does not supply a definition for arbitrary nonreductive or disconnected groups. For a finite defined Cayley degree, $\operatorname{Cay}(G)>2$ means that equivariant dominant maps exist but none has degree one or two. Empty-map classes are not assigned degree infinity without a stated convention.

The extract does not specify $k$, its characteristic, or the central-isogeny convention for simple groups. In the usual connected reductive usage, simplicity excludes nontrivial proper connected normal subgroups but can allow a finite center; no adjoint or simply connected restriction is imposed here.
\subsection{Short English statement}
Is the least degree of a conjugation-to-adjoint equivariant dominant rational map ever greater than two, and can this happen for a simple algebraic group? The field and precise simple-group convention are not supplied.
\subsection{Sources}
\begin{itemize}
\item Supplied report: \emph{Affine Algebraic Geometry} (2007), Oberwolfach Reports, DOI 10.4171/owr/2007/01, pp. 5--82; proposed year 2007. \url{https://doi.org/10.4171/owr/2007/01}.
\item Local UnsolvedMath registry, permanent ID 30000668, catalogue entry; created 25 August 2026, updated 16 September 2026. Its corrected\_verified note says both existence questions and Cayley-degree notation were preserved. The supplied 21 August 2026 literature triage labels the extracted claim open and reports no located authoritative complete resolution. No additional external reference is supplied; this dated assessment is not independently updated here.
\end{itemize}
\end{document}
